\chapter{Eight-Week Training Spine}

The training spine covers 56 active calendar days. Weekdays are full internship days. Weekend days are lighter required asynchronous work. Each day produces a concrete artifact or revision that moves the final project forward.

\section{Week 1: WiCCED Onboarding and Question Design}

\WeekHeader{1}{WiCCED onboarding and water-quality question design}{Students understand the project space, establish research-log habits, and select a feasible water-quality question.}{One-page project brief, chemical/water-quality intake table, first risk note, and Week 1 research-log entries.}{WiCCED context, water-quality variables, responsible conduct, research logs, and scope control.}

\DayPlan{1}{Orientation to WiCCED, dsDNA Core, and ONE Health Water Research}{Full day: 30 min check-in; 75 min project-space briefing; 90 min source walkthrough; 120 min research-log setup; 60 min writing; 30 min checkpoint.}{Begin the program with a concrete understanding of saltwater intrusion, Delaware water quality, coastal ecosystems, aquaculture connections, and reproducible research expectations.}{This manual, Project WiCCED website notes, research log template, project folder, and source bibliography.}{Read the program overview; build the first research-log entry; summarize why saltwater intrusion matters for water quality, aquaculture, agriculture, wetlands, and communities; list three possible student-scale questions.}{Research log opened, folder created, and three candidate questions written with one sentence of motivation each.}{The log includes date, file location, sources consulted, and at least one uncertainty that needs review.}

\DayPlan{2}{Turning Broad Interests into Testable Questions}{Full day: 30 min objective; 90 min question-design clinic; 105 min worksheet lab; 120 min independent revision; 45 min evidence writing; 30 min review.}{Convert a broad interest such as oyster health, salinity intrusion, nutrients, contaminants, or marsh water quality into a feasible 8-week research question.}{Project Brief worksheet, Water-Quality Variable Dictionary worksheet, WiCCED context chapter, and optional literature search notes.}{Draft candidate questions; identify response variables, predictors, sample units, and evidence limits; reject questions that require unavailable data or causal claims beyond the design.}{Two refined candidate questions with variables, data needs, and feasibility notes.}{Each question can be answered with available or teaching data and has a clear product by Week 8.}

\DayPlan{3}{Chemical Identity, Water-Quality Evidence, and Tentative Annotation}{Full day: 30 min check-in; 90 min concept block; 90 min uafR boundary discussion; 120 min evidence-type sorting; 60 min log update; 30 min checkpoint.}{Learn why chemical screening evidence, sensor data, nutrient measurements, and database annotations have different claim strength.}{uafR function docs, claim-template table, GC-MS annotation notes, and research log.}{Sort example claims into measured observation, tentative annotation, public database descriptor, model output, and confirmed evidence; write stronger wording for weak claims.}{Claim-evidence-limitation table for at least eight example statements.}{No statement treats a tentative match, database record, or model prediction as confirmed sample truth.}

\DayPlan{4}{Evidence Source Landscape for WiCCED Projects}{Full day: 30 min objective; 90 min source briefing; 90 min database/literature audit; 120 min project-specific source map; 60 min writing.}{Understand what PubChem, KEGG, LOTUS, FEMA, FDA/SPL, MeSH, DNREC, literature, and project datasets can and cannot support.}{Database Evidence Map worksheet, web source notes, uafR documentation, and bibliography.}{Map at least five source types to the student's project question; identify whether each source provides measurements, annotations, identifiers, pathways, occurrence, regulatory context, or literature.}{Database Evidence Map v1.}{Every source entry has a use, a limitation, and a citation or file path.}

\DayPlan{5}{Scope Gate and Project Selection}{Full day: 30 min check-in; 60 min project pitch prep; 90 min peer review; 90 min mentor review; 120 min revision; 45 min readiness checklist.}{Select one primary project question and define the minimum viable project.}{Project Brief worksheet, evidence map, risk note, file inventory, and research log.}{Present the project question; revise scope after review; define required variables, optional uafR chemical evidence, expected figures, and ML target.}{Approved Project Brief v1 and Week 2 readiness checklist.}{The project has a defined sample unit, target outcome or comparison, data source, risks, and fallback teaching-data path.}

\DayPlan{6}{Weekend Reading Annotation and Intake Revision}{Weekend: 1 to 2.5 hours.}{Consolidate Week 1 by annotating one source and revising the intake table.}{One selected reading or official project source, Chemical/Water-Quality Intake Table worksheet, and research log.}{Annotate the source; add variables, chemical names, sites, organisms, or methods that may matter; revise the intake table.}{Annotated reading note and updated intake table.}{The note separates useful facts from questions that still need verification.}

\DayPlan{7}{Weekend Reflection and Week 2 Readiness}{Weekend: 1 to 2 hours.}{Prepare to move from question design to reproducible project construction.}{Research log, Project Brief v1, evidence map, and readiness checklist.}{Write a one-page reflection on what changed during Week 1; list installation, data, and file-organization issues to solve in Week 2.}{Week 1 reflection and Week 2 action list.}{A peer or mentor can understand the project question and next steps from the project folder alone.}

\section{Week 2: R, RStudio, Git, and Reproducible Project Structure}

\WeekHeader{2}{Reproducible water-quality project build}{Students build a clean RStudio project that can be rerun from the project root.}{R project skeleton, README, script order, install/session notes, and clean-session rerun record.}{RStudio, base R, uafR readiness, file paths, script hygiene, Git habits, and bundle use.}

\DayPlan{8}{R/RStudio and uafR Readiness Check}{Full day: 30 min check-in; 90 min installation audit; 90 min setup script; 120 min project skeleton; 60 min documentation.}{Make sure the student's local machine can run the training scripts and find uafR.}{\file{scripts/00_setup_check.R}, \file{scripts/01_project_setup.R}, RStudio, student bundle, and project folder.}{Run the setup check; create the project skeleton; record R version, package status, and any missing dependencies.}{Project folder with subfolders, setup log, and session note.}{A clean R session can run the setup check and report uafR status without hidden local paths.}

\DayPlan{9}{Reading, Inspecting, Cleaning, and Writing Tables in R}{Full day: 30 min objective; 90 min table concepts; 120 min guided lab; 120 min independent practice; 45 min log.}{Learn how to inspect CSV tables before using them in analysis or models.}{Teaching water-quality CSV, \file{scripts/02_water_quality_qaqc.R}, data dictionary worksheet, and RStudio.}{Read the teaching dataset; inspect dimensions, column names, types, missingness, ranges, and units; write a cleaned copy and QA/QC summary.}{Cleaned teaching-data copy and first QA/QC summary CSV.}{The student can explain each column, its unit, its missingness, and one suspicious-value rule.}

\DayPlan{10}{Git Basics, Safe Paths, and Repository Habits}{Full day: 30 min check-in; 75 min Git concept block; 90 min practical lab; 120 min README revision; 60 min peer handoff.}{Use version control habits that protect data, scripts, and reproducibility.}{Git or RStudio Git pane, README draft, .gitignore template, and research log.}{Initialize or inspect a local repository; make a non-sensitive commit if Git is available; update README with script order and data rules; document what files should never be committed.}{README v1, .gitignore note, and commit or Git-status screenshot/note.}{No private raw data, credentials, machine-specific cache folders, or absolute personal paths are required to rerun the project.}

\DayPlan{11}{Numbered Workflows and Raw Versus Processed Data}{Full day: 30 min objective; 90 min workflow design; 120 min script-order lab; 90 min independent organization; 60 min writing.}{Create a script order that preserves raw data and makes processed outputs traceable.}{Project folder, README, teaching data, and script templates.}{Define \file{data_raw}, \file{data_processed}, \file{results}, \file{figures}, \file{models}, \file{logs}, and \file{reports}; add a script index and expected outputs.}{Script-order table in README and updated folder tree.}{A new user can tell which script to run first and which files are inputs versus outputs.}

\DayPlan{12}{Clean-Session Reproducibility Check}{Full day: 30 min check-in; 90 min rerun protocol; 120 min clean-session rerun; 120 min debugging; 45 min sign-off note.}{Verify that Week 2 work can be rerun without relying on hidden console objects.}{RStudio restart, setup scripts, QA/QC script, README, and research log.}{Restart R; run scripts in order; fix path errors; record failures and corrections; produce a clean-session rerun note.}{Clean-session rerun record and reproducibility checkpoint.}{All Week 2 outputs can be recreated from scripts and documented inputs.}

\DayPlan{13}{Weekend Debugging Log and Session Information}{Weekend: 1 to 2 hours.}{Turn setup issues into documented knowledge.}{Research log, session information, console notes, README.}{Record R version, package versions, operating system, unresolved errors, and command recap.}{Weekend debugging log and session-info note.}{The log is specific enough that another student can reproduce the issue or verify that it is fixed.}

\DayPlan{14}{Weekend Peer Handoff}{Weekend: 1 to 2.5 hours.}{Test whether the project folder is understandable without verbal explanation.}{README, folder tree, setup log, and peer review form.}{Ask a peer or mentor to inspect the folder from README alone; collect questions; revise confusing sections.}{Peer handoff note and README revision.}{The reviewer can identify the project question, data location, script order, and current outputs.}

\section{Week 3: Water-Quality Data Literacy and QA/QC}

\WeekHeader{3}{Water-quality data literacy and QA/QC}{Students build a defensible water-quality dataset and understand field/lab metadata.}{Water-quality variable dictionary, QA/QC flag log, missingness report, and data-readiness gate.}{Units, detection limits, field metadata, method effects, missingness, outliers, and Delaware monitoring context.}

\DayPlan{15}{Sample Units, Sites, Dates, and Metadata}{Full day: 30 min check-in; 90 min metadata concept block; 120 min data dictionary lab; 120 min project update; 45 min writing.}{Define what one row means and which metadata are required for interpretation.}{Teaching data, project data if available, Variable Dictionary worksheet, and QA/QC Flag Log.}{Identify sample unit, site, date/time, depth, tide/rainfall context if available, method, units, and responsible data steward.}{Water-Quality Variable Dictionary v1.}{Every variable has a unit, type, source, method note, and interpretation caveat.}

\DayPlan{16}{Ranges, Units, Detection Limits, and Suspicious Values}{Full day: 30 min objective; 90 min QA/QC rules; 120 min script lab; 120 min project-specific flags; 60 min review.}{Create transparent rules for impossible, suspicious, missing, and censored water-quality values.}{\file{scripts/02_water_quality_qaqc.R}, QA/QC worksheet, teaching data, and project data if approved.}{Run QA/QC checks; adapt range rules cautiously; flag values without deleting them; write a decision note for each rule.}{QA/QC flag log and flagged-data output.}{The student can explain why each flag exists and can recover the original raw value.}

\DayPlan{17}{Salinity, Conductivity, Dissolved Oxygen, and Nutrients}{Full day: 30 min check-in; 90 min variable interpretation; 90 min exploratory lab; 120 min independent summaries; 60 min writing.}{Understand core variables often used in saltwater intrusion and aquatic ecosystem interpretation.}{Teaching data, DNREC monitoring context notes, variable dictionary, and plotting script template.}{Summarize salinity, conductivity, dissolved oxygen, pH, turbidity, nitrate, chlorophyll, and temperature by site/date/context; identify likely confounders.}{Exploratory summary table and two draft plots.}{No plot or summary implies causation; all units and sample counts are visible.}

\DayPlan{18}{Missingness, Bias, and Sampling Design}{Full day: 30 min objective; 90 min missingness concepts; 120 min missingness audit; 120 min project risk revision; 45 min checkpoint.}{Learn how missing values, uneven sampling, season, site, and method can bias analysis and ML models.}{Source Coverage/Missingness worksheet, teaching data, project data, and research log.}{Compute missingness by variable, site, and date; identify structural versus accidental missingness; decide which variables are usable.}{Missingness report and project risk update.}{Variables with heavy or biased missingness are either excluded or clearly marked as exploratory.}

\DayPlan{19}{Data-Readiness Gate}{Full day: 30 min check-in; 90 min review preparation; 90 min mentor review; 120 min revision; 60 min documentation.}{Freeze the first analysis-ready water-quality dataset or approve the teaching-data fallback.}{Variable dictionary, QA/QC log, missingness report, README, and cleaned data.}{Present the data pipeline; answer what each row means; revise after review; define the Week 4 chemical/uafR path.}{Data-readiness gate memo.}{The memo states which dataset will be used, which variables are included/excluded, and why.}

\DayPlan{20}{Weekend Reading Annotation on Water Quality}{Weekend: 1 to 3 hours.}{Connect data literacy to a source on Delaware water monitoring, aquaculture, marsh systems, salinity, or machine learning.}{Approved reading, research log, and evidence map.}{Annotate one source; extract variables, methods, and limitations; connect it to the project dictionary.}{Annotated reading and updated dictionary notes.}{The annotation records how the source changes variable interpretation or project scope.}

\DayPlan{21}{Weekend Data Dictionary Revision}{Weekend: 1 to 2 hours.}{Clean up Week 3 artifacts before chemical screening or modeling begins.}{Variable dictionary, QA/QC flag log, missingness report, README.}{Resolve inconsistent variable names; revise units; add missing caveats; prepare Week 4 questions.}{Data dictionary v2 and Week 4 readiness checklist.}{A reader can understand the analysis-ready data without opening the raw file first.}

\section{Week 4: uafR Chemical Screening and Evidence Mapping}

\WeekHeader{4}{uafR chemical evidence for water-quality questions}{Students learn when uafR is relevant, run an offline workflow, and create a cautious chemical evidence map.}{uafR workflow output or decision note, chemical evidence map, query settings log, and claim-limit memo.}{GC-MS hit tables, tentative annotation, PubChem profiles, categorate research output, validation, and trait evidence.}

\DayPlan{22}{GC-MS Hit Tables and Water-Quality Chemical Screening}{Full day: 30 min objective; 90 min GC-MS hit-table concept block; 120 min input audit lab; 120 min project relevance decision; 45 min writing.}{Understand what GC-MS hit tables can add to a water-quality project and when they are not needed.}{uafR docs, GC-MS Input Column Audit worksheet, optional project chemical data, and teaching examples.}{Inspect required columns for \code{spreadOut()}; identify compound names, retention time, area, match factor, file names, and suspicious rows; decide whether chemical screening belongs in the project.}{GC-MS input audit or chemical-screening not-applicable note.}{The project does not force uafR chemical analysis if no defensible chemical input exists.}

\DayPlan{23}{Simulated Core uafR-Compatible Workflow}{Full day: 30 min check-in; 90 min function map; 120 min guided run; 120 min object inspection; 60 min log.}{Run and inspect a simulated WiCCED-themed GC-MS workflow without relying on live web services.}{\file{data/wicced_simulated_gcms_hits.csv}, simulated query/reference files, \code{spreadOut()}, \code{mzExacto()}, and \file{scripts/03_uafr_chemical_screening_template.R}.}{Run the simulated pipeline; inspect the raw hit table, spread RDS, exact-match output, long abundance table, context summary, and evidence notes; record which operations mirror uafR hit-table processing and which live functions require web lookup.}{Simulated GC-MS workflow note and saved output tables.}{The student can explain each required input column, why duplicate detections are aggregated, and why simulated chemical patterns are not environmental evidence.}

\DayPlan{24}{Research-Grade Enrichment and Validation First}{Full day: 30 min objective; 90 min enrichment concept block; 90 min small-query planning; 120 min optional live/saved result inspection; 60 min validation notes.}{Learn the safe pattern for \code{categorate(detail = "research")} and validation outputs.}{uafR docs, \code{validateCategorateResult()}, Query Settings Log, saved result if available, and cache/throttle notes.}{Plan a tiny query list; use saved output or run only if network and package dependencies are ready; inspect \code{ValidationSummary}, \code{SourceCoverage}, \code{TableQuality}, \code{DataDictionary}, \code{ChemicalTraits}, and \code{ChemicalTraitEvidence}.}{Categorate validation note and query settings log.}{No live web query is scaled until the small run or saved example is understood and documented.}

\DayPlan{25}{Chemical Traits for Water-Quality Grouping}{Full day: 30 min check-in; 90 min trait concepts; 120 min matrix lab; 120 min project-specific evidence mapping; 45 min review.}{Convert source-backed chemical annotations into variables that could be used for grouping or modeling.}{\code{chemicalTraitMatrix()}, \code{chemicalTraitEvidence()}, Trait Evidence Audit worksheet, and saved categorate result if available.}{Inspect trait names, confidence, sources, evidence text, and missingness; build or inspect a binary/count/confidence trait matrix; reject narrative-only fields as model features unless discretized reproducibly.}{Trait Evidence Audit v1.}{Every retained trait has a source table, source field, extraction rule, confidence note, and caveat.}

\DayPlan{26}{Claim-Limit Memo for Chemical Evidence}{Full day: 30 min objective; 90 min strong/weak claims; 90 min evidence court; 120 min writing; 60 min checkpoint.}{Write a project-specific memo that prevents overclaiming from tentative chemical evidence.}{Evidence Court worksheet, claim templates, uafR outputs or decision note, and research log.}{Trace five possible chemical statements back to evidence; rewrite weak claims; list confirmation work required for stronger claims.}{Chemical claim-limit memo.}{The memo clearly distinguishes tentative matches, database descriptors, source-backed traits, and confirmed identities.}

\DayPlan{27}{Weekend Wrong-Input Troubleshooting}{Weekend: 1 to 2.5 hours.}{Practice diagnosing malformed chemical-screening inputs without damaging raw data.}{Wrong-Column GC-MS Input Repair lab, example malformed table, and research log.}{Identify missing or mislabeled columns; write what the error means; create a corrected copy with documented changes.}{Troubleshooting note and corrected example file.}{The raw file remains unchanged and every correction is documented.}

\DayPlan{28}{Weekend Chemical Evidence Archive}{Weekend: 1 to 2 hours.}{Prepare chemical evidence for modeling decisions in Week 5.}{uafR notes, query settings log, evidence map, README.}{Archive saved outputs or write the not-applicable note; update README with chemical-evidence status.}{Chemical evidence archive and Week 5 readiness checklist.}{The project folder states exactly whether uafR chemical evidence will enter the ML feature table.}

\section{Week 5: Machine-Learning Foundations for Water Quality}

\WeekHeader{5}{Machine-learning foundations}{Students build a baseline model from a documented feature table and evaluate it honestly.}{ML feature table v1, train/test split plan, baseline regression or classification model, and first model card.}{Prediction targets, leakage, baseline models, metrics, cross-validation concepts, and model cards.}

\DayPlan{29}{Prediction Questions and Target Definition}{Full day: 30 min check-in; 90 min ML framing; 120 min target-design lab; 120 min feature inventory; 45 min writing.}{Define a prediction or classification target that is meaningful and available before interpretation.}{ML Feature Inventory worksheet, Train/Test Split Plan, teaching data, project data, Week 3 outputs, and \file{ml_resources/ml_glossary.csv}.}{Choose one target such as salinity, dissolved oxygen, nitrate, or high-salinity class; list candidate predictors; remove features that duplicate or reveal the target; read the glossary terms for target, feature, split, baseline, leakage, and model card.}{ML Feature Inventory v1 and target definition note.}{The target is measurable, not circular, and not created from predictors that will also be used as model inputs.}

\DayPlan{30}{Feature Table and Split Script}{Full day: 30 min objective; 90 min leakage examples; 120 min guided script run; 120 min output inspection; 60 min review.}{Protect model evaluation from overly optimistic results by building a documented feature table and split plan.}{Train/Test Split Plan worksheet, teaching data, site/date metadata, \file{scripts/04_ml_01_build_feature_table.R}, and \file{ml_resources/ml_workflow_checklist.csv}.}{Run the feature-table script; open the feature dictionary, training rows, testing rows, and split plan; identify leakage risks from site labels, dates, conductivity, post-outcome variables, duplicates, and repeated sampling.}{Feature table, feature dictionary, and approved train/test split plan.}{The split matches the intended use case and every included/excluded variable has a documented reason.}

\DayPlan{31}{Explore Before Modeling}{Full day: 30 min check-in; 90 min exploration concepts; 120 min guided script run; 120 min plot/table inspection; 45 min log.}{Inspect the feature table before fitting models.}{\file{scripts/04_ml_02_explore_features.R}, feature table, feature dictionary, exploration plots, and research log.}{Run the exploration script; inspect feature summaries, correlations, class balance, and plots; answer the beginner exploration questions before any model interpretation.}{Feature summary, correlation table, class balance table, exploration plots, and answered exploration questions.}{The student can name at least one plausible feature, one weak feature, one excluded feature, and one limitation before modeling.}

\DayPlan{32}{Train Regression and Classification Models}{Full day: 30 min objective; 90 min regression/classification concepts; 120 min guided model run; 120 min metrics review; 60 min writing.}{Build simple baseline models and compare them to candidate models.}{\file{scripts/04_ml_03_train_models.R}, feature table, split plan, and glossary.}{Run the training script; compare mean-only regression to linear regression; compare majority-class classification to logistic classification; inspect RMSE, MAE, confusion-matrix counts, sensitivity, specificity, and balanced accuracy.}{Regression metrics, classification metrics, predictions table, coefficients table, and model-comparison summary.}{The student can explain false positives, false negatives, baseline comparison, and why test-set performance matters.}

\DayPlan{33}{Diagnostics, Model Card, and First ML Review}{Full day: 30 min check-in; 90 min diagnostic concepts; 90 min script run; 120 min model-card review; 60 min checkpoint.}{Document the first model so its purpose, data, limits, and performance are visible.}{\file{scripts/04_ml_04_diagnostics_and_model_card.R}, \file{ml_resources/model_card_field_guide.md}, metrics CSVs, plots, predictions, and research log.}{Run the diagnostics script; inspect observed-versus-predicted plots, residuals, classification probabilities, diagnostic summary, model card, and interpretation guide; revise features, split, metrics, and limitations; list next experiments for Week 6.}{Diagnostic plots, diagnostic summary, Model Card v1, and ML review note.}{The model card names the target, rows used, exclusions, features, split design, metrics, diagnostics, limitations, and prohibited uses.}

\DayPlan{34}{Weekend ML Reading Annotation}{Weekend: 1 to 3 hours.}{Connect the baseline model to water-quality ML literature without copying claims.}{Approved ML review or methods article, model card, research log.}{Annotate one ML water-quality source; identify target types, algorithms, validation designs, and uncertainty concerns.}{Annotated ML reading and model-card update.}{The annotation improves the student's validation plan or feature rationale.}

\DayPlan{35}{Weekend ML Rerun and Sensitivity Note}{Weekend: 1 to 2.5 hours.}{Check whether the model result changes under a simple defensible sensitivity run.}{Beginner ML scripts, feature inventory, model card, and research log.}{Rerun the ladder after removing one feature or changing one documented setting; record differences without cherry-picking.}{Sensitivity note.}{The note reports what changed, whether the baseline comparison changed, and why the preferred result should remain cautious.}

\section{Week 6: Feature Engineering, Model Comparison, and Spatial/Temporal Context}

\WeekHeader{6}{Feature engineering and model comparison}{Students build an analysis-ready feature table, compare models, and inspect spatial/temporal context without leakage.}{Frozen feature table, model comparison table, leakage audit, and model interpretation memo.}{Feature engineering, grouping variables, chemical traits, spatial context, temporal structure, model comparison, and uncertainty.}

\DayPlan{36}{Feature Table Build}{Full day: 30 min objective; 90 min feature-engineering rules; 120 min table build; 120 min independent revision; 60 min review.}{Create one feature table used by all Week 6 modeling.}{ML Feature Inventory, water-quality cleaned data, optional uafR trait matrix, and README.}{Join only approved features; document transformations; create missingness indicators only when justified; save the feature table with a data dictionary.}{Feature table v1 and feature dictionary.}{Each feature has a source, unit or encoding, transformation rule, and leakage status.}

\DayPlan{37}{Chemical Traits as Optional ML Features}{Full day: 30 min check-in; 90 min trait-to-feature caution; 120 min feature integration lab; 120 min evidence audit; 45 min log.}{Use uafR-derived features only when the chemical evidence and sample mapping are defensible.}{Saved uafR trait matrix if available, Trait Evidence Audit, feature table, and chemical claim-limit memo.}{Decide whether chemical traits enter the model; if they do, join by documented sample or compound mapping and preserve missingness; if not, explain why.}{Chemical feature decision note.}{No compound trait enters the model without a traceable link to sample-level evidence and source-backed trait evidence.}

\DayPlan{38}{Advanced Validation After the Beginner Model}{Full day: 30 min objective; 90 min model comparison concepts; 120 min advanced validation run; 120 min diagnostics; 60 min writing.}{Compare the beginner model against stricter validation checks without losing interpretability.}{Beginner ML outputs, \file{scripts/05_ml_01_advanced_validation.R}, \file{ml_resources/advanced_ml_extension_guide.md}, model-card template, and metrics table.}{Confirm that the beginner ladder is understood; run the advanced validation script; compare time-holdout and leave-one-site-out results; record where blocked validation weakens or supports the beginner result.}{Advanced validation folds, split comparison table, and validation notes.}{All comparisons use documented targets, eligible rows, and feature sets; any skipped model is treated as evidence of overfitting risk rather than a failure to hide.}

\DayPlan{39}{Spatial and Temporal Context}{Full day: 30 min check-in; 90 min spatial/temporal concepts; 120 min site/date summaries; 120 min model-risk review; 45 min log.}{Understand how site, season, tide, rainfall, and repeated sampling can structure model results.}{Site/date metadata, feature table, map or site-context notes, and leakage audit.}{Summarize target and predictors by site/date/context; identify whether site identifiers act as shortcuts; decide whether site/date features belong in the main model or sensitivity analysis.}{Spatial/temporal context memo.}{The memo distinguishes environmental context from leakage-prone identifiers.}

\DayPlan{40}{Sensitivity, Permutation Importance, and Leakage Review}{Full day: 30 min objective; 90 min leakage audit; 90 min guided sensitivity run; 120 min output review; 60 min sign-off.}{Test whether feature-set changes improve performance for defensible reasons or because they leak information.}{Feature table, split plan, advanced sensitivity script, leakage checklist, and research log.}{Run the sensitivity script; compare the core feature set with conductivity and context sensitivity outputs; inspect permutation importance; write why a high-importance feature is not causal proof.}{Feature-set sensitivity table, permutation-importance table, sensitivity warnings, and leakage-audit memo.}{Conductivity or context features are not promoted to the main model unless the project question justifies their availability and interpretation.}

\DayPlan{41}{Weekend Matrix Sparsity and Missing-Data Challenge}{Weekend: 1 to 3 hours.}{Evaluate whether sparse features are helping or just adding noise.}{Feature table, missingness report, optional trait matrix, and model card.}{Calculate missingness or prevalence for each feature; identify sparse features; propose exclusions or sensitivity checks.}{Sparsity and missingness sensitivity note.}{Any proposed feature removal is justified before seeing whether it improves the preferred result.}

\DayPlan{42}{Weekend Advanced Model Review}{Weekend: 1 to 2.5 hours.}{Translate Week 6 results into cautious interpretation.}{Advanced model outputs, \file{scripts/05_ml_03_advanced_model_report.R}, decision-gate table, leakage audit, and research log.}{Run the advanced report script; open the model review and decision gate; write a one-page memo naming the most defensible model, why it was chosen, and what it cannot support.}{Advanced model review, decision gate, and model interpretation memo.}{The memo discusses uncertainty, leakage, split design, and limitations at least as clearly as performance.}

\section{Week 7: Visualization, Evidence Synthesis, and Communication}

\WeekHeader{7}{Visualization and evidence synthesis}{Students turn processed data and model outputs into figures, tables, and defensible scientific interpretation.}{Figure package, claim-evidence-limitation memo, literature synthesis table, and draft report sections.}{Visual honesty, missingness, model diagnostics, source evidence, literature integration, and reviewer-style critique.}

\DayPlan{43}{Figure Planning Before Plotting}{Full day: 30 min check-in; 90 min visual-purpose clinic; 120 min figure planning; 120 min independent sketching; 45 min review.}{Plan figures around claims, not around available plotting options.}{Figure Planning Worksheet, final feature table, model outputs, and source evidence.}{Define the question for each figure; identify audience, variables, grouping, missingness, caption claim, and caveat.}{Figure plan for at least four candidate figures.}{Each planned figure has a purpose, data source, caption draft, and claim limit.}

\DayPlan{44}{Water-Quality and Source-Coverage Figures}{Full day: 30 min objective; 90 min plotting principles; 120 min guided plotting; 120 min project plotting; 60 min log.}{Create figures that show water-quality patterns and data coverage honestly.}{Cleaned data, missingness report, source coverage table if available, and plotting script.}{Build summaries by site/date/context; plot salinity/conductivity/DO/nutrients with units; create a missingness or source-coverage heatmap.}{Two draft water-quality/source-coverage figures.}{Axes, units, sample counts, missingness, and grouping definitions are visible.}

\DayPlan{45}{Model Diagnostic and Interpretation Figures}{Full day: 30 min check-in; 90 min diagnostics; 120 min model plot lab; 120 min independent revision; 45 min checkpoint.}{Visualize model performance without exaggerating accuracy.}{Model metrics, predictions, model card, and figure caption checklist.}{Plot observed versus predicted values, residuals, confusion matrix, or probability distributions as appropriate; write captions that state split design and limits.}{Two draft model figures and caption drafts.}{No figure reports training-only performance as final evidence.}

\DayPlan{46}{Figure Clinic and Peer Review}{Full day: 30 min objective; 90 min peer critique; 90 min revision; 120 min caption repair; 60 min checkpoint.}{Improve figures through reviewer-style feedback.}{Figure Caption Checklist, draft figures, peer review form, and claim memo.}{Review axes, legends, color choices, caption claims, missingness, and accessibility; revise figures and captions.}{Revised figure package v1.}{A reviewer can identify the data, variables, sample counts, and claim limit from the figure and caption.}

\DayPlan{47}{Literature-to-Evidence Bridge}{Full day: 30 min check-in; 90 min literature strategy; 120 min source extraction; 120 min writing; 45 min review.}{Connect project results to literature without overstating similarity.}{Literature and Evidence Synthesis worksheet, annotated readings, project figures, and model card.}{Search by variables, site context, organism/system, method, and model target; extract system, method, chemical/water evidence, main claim, and limitation.}{Literature synthesis table and one claim-evidence-limitation paragraph.}{The paragraph cites evidence and states how the student's dataset is similar or different.}

\DayPlan{48}{Weekend Figure Revision and Export}{Weekend: 1 to 3 hours.}{Prepare figures for final report drafting.}{Draft figures, captions, README, and figure export settings.}{Export figures with consistent names; update captions; record script and output paths in README.}{Figure package v2.}{Each figure can be regenerated and has a matching caption.}

\DayPlan{49}{Weekend Visual Narrative Outline}{Weekend: 1 to 2 hours.}{Turn figures into the skeleton of the final report.}{Figure package, literature synthesis table, model interpretation memo.}{Order figures; write one sentence per figure explaining what it contributes to the story.}{Visual narrative outline.}{The outline does not contain unsupported causal or confirmation claims.}

\section{Week 8: Capstone Report, Presentation, and Handoff}

\WeekHeader{8}{Final research product and handoff}{Students finish a reproducible research package that another researcher can inspect and continue.}{Manuscript-style report, presentation/poster, reproducible folder, final package audit, continuation note, and exit reflection.}{Scientific writing, reproducibility audit, presentation, peer review, archive discipline, and next-step planning.}

\DayPlan{50}{Report Storyboard}{Full day: 30 min objective; 90 min report structure; 120 min storyboard lab; 120 min drafting; 45 min checkpoint.}{Organize the final report before writing full prose.}{Report template, visual narrative outline, project brief, figures, and model card.}{Draft title, abstract bullets, introduction rationale, methods sections, results sequence, discussion claims, limitations, and reproducibility statement.}{Report storyboard.}{Every planned result connects to a figure/table and every claim connects to evidence.}

\DayPlan{51}{Methods From Scripts and Logs}{Full day: 30 min check-in; 90 min methods standards; 120 min script-to-methods lab; 120 min writing; 60 min review.}{Write methods that are specific enough to rerun.}{README, scripts, data dictionary, QA/QC log, uafR query settings, model card, and session info.}{Write data sources, cleaning rules, uafR settings if used, feature engineering, train/test split, model type, metrics, and software versions.}{Methods section draft.}{A reader can match each methods paragraph to a script, file, or documented setting.}

\DayPlan{52}{Results Without Overclaiming}{Full day: 30 min objective; 90 min results-writing clinic; 120 min paragraph drafting; 120 min table/figure checks; 45 min log.}{Write results that describe observations and model outputs without causal inflation.}{Figures, tables, model metrics, claim-evidence-limitation memo, and caption checklist.}{Write results paragraphs tied to figure numbers; include sample sizes, units, model split, and uncertainty; remove unsupported explanations.}{Results section draft.}{Each result paragraph describes what was observed or predicted and avoids unsupported mechanism claims.}

\DayPlan{53}{Discussion, Limitations, and Next Steps}{Full day: 30 min check-in; 90 min discussion clinic; 120 min writing; 90 min peer critique; 60 min revision.}{Explain why the results matter while keeping uncertainty visible.}{Literature synthesis table, model interpretation memo, chemical claim-limit memo, and project risk note.}{Write discussion paragraphs that connect to WiCCED context; include limitations for data, chemistry, model validation, and generalization; define follow-up tests.}{Discussion and limitations draft.}{Limitations are concrete and tied to the actual dataset and methods.}

\DayPlan{54}{Final Reproducibility Audit}{Full day: 30 min objective; 90 min audit protocol; 150 min clean-session rerun; 90 min fixes; 45 min sign-off.}{Verify the project can be rerun from a clean session.}{\file{scripts/run_wicced_smoke_tests.R}, project scripts, README, final package audit worksheet, and session info.}{Restart R; rerun scripts; compare expected outputs; fix path or missing-file problems; document remaining risks.}{Final reproducibility audit and session information.}{All required outputs exist or a clear not-applicable note explains why.}

\DayPlan{55}{Presentation or Poster Build}{Full day: 30 min check-in; 90 min presentation structure; 120 min slide/poster build; 90 min practice; 60 min revision.}{Create a clear oral or poster presentation around one central result.}{Report draft, figure package, model card, and presentation planning worksheet.}{Build slides/poster with question, data, methods, main result, model validation, limitations, and next steps; practice a 5 to 8 minute version.}{Presentation/poster draft and speaking notes.}{The presentation makes one central argument and includes limitations without hiding them.}

\DayPlan{56}{Final Presentation, Archive, and Continuation Note}{Full day: 30 min setup; 90 min presentations; 90 min Q\&A and feedback; 120 min final archive; 60 min exit reflection.}{Complete the program with a usable research package and a clear path for continuation.}{Final report, presentation, project folder, final package audit, continuation note, and research log.}{Present the project; answer reviewer-style questions; archive the package; write unresolved issues and recommended next experiments.}{Final report, presentation, archive, continuation note, and exit reflection.}{A future student can open the folder, understand the project, rerun the scripts, and see exactly what remains unresolved.}
